% =====================================================================
%  BUỔI 13 — KẾ HOẠCH GIẢNG DẠY (từ giao-vien/01-ke-hoach-chua-de.md)
%  Sinh tự động bằng scripts/md2tex-chung.py — không sửa tay.
% =====================================================================
\documentclass[11pt, a4paper]{article}
% =====================================================================
%  CAE SHANGHAI — GÓI LỆNH DÙNG CHUNG CHO TÀI LIỆU LUYỆN THI CSCA TOÁN
%  Biên dịch bằng XeLaTeX. Mọi file .tex trong buổi học \input file này.
% =====================================================================

% ---------- NGÔN NGỮ & PHÔNG CHỮ ----------
\usepackage{fontspec}
\usepackage{polyglossia}
\setmainlanguage{vietnamese}
\setotherlanguage{english}

% Chữ Latin (Anh + Việt): TeX Gyre Termes — bản clone Times New Roman,
% đi kèm TeX Live, hỗ trợ đầy đủ dấu tiếng Việt.
\setmainfont{TeX Gyre Termes}
\setsansfont{TeX Gyre Heros}

% Chữ Hán: ưu tiên Source Han Serif SC (思源宋体) — chuẩn học thuật Trung Quốc.
% Nếu máy chưa cài, tự động dùng Songti SC (宋体, có sẵn trên macOS).
% Cả hai đều thuộc họ 宋体 nên diện mạo tài liệu không đổi.
\IfFontExistsTF{Source Han Serif SC}{
  \newcommand{\cjkfamilyname}{Source Han Serif SC}
}{
  \IfFontExistsTF{Songti SC}{
    \newcommand{\cjkfamilyname}{Songti SC}
  }{
    \newcommand{\cjkfamilyname}{Noto Serif CJK SC}
  }
}

% xeCJK tự chuyển font khi gặp ký tự Hán, không cần bọc \cjk{} thủ công.
% Đây là điều kiện bắt buộc: font Latin (TeX Gyre Termes) không có glyph chữ Hán,
% nếu không khai báo xeCJK thì mọi chữ Hán viết trần sẽ bị nuốt khỏi bản in.
\usepackage{xeCJK}
\setCJKfamilyfont{caehan}{\cjkfamilyname}[Scale=0.95]
\setCJKmainfont{\cjkfamilyname}[Scale=0.95]
\setCJKsansfont{\cjkfamilyname}[Scale=0.95]
\setCJKmonofont{\cjkfamilyname}[Scale=0.95]

% \cjk{} ép dùng font Hán trong ngữ cảnh cần thiết (ví dụ trong tiêu đề in đậm
% hoặc trong bảng có định dạng riêng). Dùng \CJKfamily chứ không phải tên font,
% vì tên font chỉ là văn bản — sẽ in ra chữ thay vì đổi font.
\newcommand{\cjk}[1]{{\CJKfamily{caehan}#1}}

% ---------- TOÁN HỌC ----------
\usepackage{amsmath, amssymb, amsthm, mathtools}
\usepackage{unicode-math}
\setmathfont{TeX Gyre Termes Math}

% ---------- BỐ CỤC TRANG ----------
% Lề phải chừa vùng an toàn của letterhead:
%   banner trên chiếm 0–19mm  → top tối thiểu 2,2cm
%   banner dưới chiếm 281–297mm → bottom tối thiểu 2,2cm
% Lề trên 2,6cm (26mm) để tiêu đề đầu trang không đè banner (banner hết ở 19mm).
\usepackage[a4paper, top=2.6cm, bottom=2.4cm, left=2.5cm, right=2.0cm]{geometry}
\usepackage{fancyhdr}
\usepackage{titlesec}
\usepackage{enumitem}
\usepackage{graphicx}
\usepackage{booktabs, array, longtable, tabularx}
\usepackage{xcolor}
\usepackage[most]{tcolorbox}
\usepackage[colorlinks=true, linkcolor=cae_blue, citecolor=cae_blue, urlcolor=cae_blue]{hyperref}

% ---------- BẢNG MÀU THƯƠNG HIỆU ----------
\definecolor{cae_blue}{RGB}{0, 51, 102}
\definecolor{cae_gold}{RGB}{196, 149, 60}
\definecolor{box_bg}{RGB}{245, 247, 250}
\definecolor{box_line}{RGB}{0, 51, 102}

% ---------- TIÊU ĐỀ MỤC ----------
% Đệm đầu trang: banner letterhead chiếm 0–19mm, mà lề trên là 22mm.
% \section* đầu trang (sau \clearpage) sát lề nên phải chừa thêm, nếu không
% sẽ đè lên banner.
\titleformat{\section}
  {\normalfont\large\bfseries\sffamily\color{cae_blue}}{\thesection.}{0.6em}{}
\titleformat{\subsection}
  {\normalfont\normalsize\bfseries\sffamily\color{cae_blue!85!black}}{\thesubsection}{0.6em}{}
\titlespacing*{\section}{0pt}{14pt}{6pt}
\titlespacing*{\subsection}{0pt}{10pt}{4pt}

% Tiêu đề không đánh số ở đầu trang: \section* sau \clearpage nằm sát lề trên
% nên đè lên banner letterhead. Dùng \vspace* ngay TRƯỚC \section* không đủ,
% vì titlesec chèn khoảng cách riêng. Cách chắc chắn là đệm bằng \null.
\newcommand{\sectiondautrang}[1]{\section*{#1}}

% Dạng bài dùng \subsection*{} nên cần định dạng riêng, đậm và có gạch chân màu
\titleformat{\subsubsection}
  {\normalfont\normalsize\bfseries\sffamily\color{cae_blue!85!black}}{\thesubsubsection}{0.6em}{}

% ---------- LETTERHEAD (ảnh nền toàn trang) ----------
% Ảnh letterhead của CAE Shanghai: banner xanh trên/dưới kèm logo chìm giữa.
% Tỉ lệ ảnh 1414x2000 = 0.7070, khớp A4 (0.7071) nên chèn thẳng không méo.
% Đặt ở lớp nền dưới cùng để chữ không bị ảnh che.
% Dùng \AddToShipoutPictureBG (không có dấu *) để áp dụng cho MỌI trang.
\usepackage{eso-pic}
\AddToShipoutPictureBG{%
  \AtPageLowerLeft{%
    \includegraphics[width=\paperwidth,height=\paperheight]{images/letterhead.png}%
  }%
}

% ---------- ĐẦU / CHÂN TRANG ----------
% Letterhead đã có banner trên và dưới, nên bỏ dòng chạy để tránh chồng chữ.
% Số trang đặt trong vùng nội dung, sát trên banner dưới, lề phải.
\pagestyle{fancy}
\setlength{\headheight}{0pt}
\fancyhf{}
\renewcommand{\headrulewidth}{0pt}
\renewcommand{\footrulewidth}{0pt}
\fancyfoot[R]{\raisebox{0.9cm}{\small\color{cae_blue}\thepage}}

% ---------- MÔI TRƯỜNG NỘI DUNG ----------
% Trong tcolorbox, LaTeX đặt parindent = parskip = 0, nên các đoạn văn trong
% khung dính liền nhau. Đặt lại khoảng cách đoạn cho mọi khung.
\newcommand{\khoangcachdoan}{\setlength{\parskip}{0.5em}\setlength{\parindent}{0pt}}
\tcbset{before upper={\khoangcachdoan}}

\newtcolorbox{dinhnghia}[1][]{
  colback=box_bg, colframe=box_line, boxrule=0.6pt, arc=1pt,
  fonttitle=\bfseries, title={Định nghĩa #1}, breakable}
\newtcolorbox{dinhly}[1][]{
  colback=box_bg, colframe=cae_gold, boxrule=0.6pt, arc=1pt,
  fonttitle=\bfseries, title={Tính chất #1}, breakable}
\newtcolorbox{vidu}[1][]{
  colback=white, colframe=black!40, boxrule=0.4pt, arc=1pt,
  fonttitle=\bfseries, title={Ví dụ #1}, breakable}
\newtcolorbox{baitap}[1][]{
  colback=white, colframe=cae_blue!50, boxrule=0.4pt, arc=1pt,
  fonttitle=\bfseries, title={Câu #1}, breakable}

% Ô đáp án: viền xanh đậm hơn, dùng trong file đáp án
\newtcolorbox{dapan}[1][]{
  colback=box_bg, colframe=cae_blue, boxrule=0.6pt, arc=1pt,
  fonttitle=\bfseries, title={Câu #1}, breakable}

% ---------- TÊN CHỨNG MINH ----------
\renewcommand{\proofname}{\normalfont\bfseries Chứng minh}

% ---------- MÔN HỌC ----------
% Đổi được ở file .tex trước khi % =====================================================================
%  CAE SHANGHAI — GÓI LỆNH DÙNG CHUNG CHO TÀI LIỆU LUYỆN THI CSCA TOÁN
%  Biên dịch bằng XeLaTeX. Mọi file .tex trong buổi học \input file này.
% =====================================================================

% ---------- NGÔN NGỮ & PHÔNG CHỮ ----------
\usepackage{fontspec}
\usepackage{polyglossia}
\setmainlanguage{vietnamese}
\setotherlanguage{english}

% Chữ Latin (Anh + Việt): TeX Gyre Termes — bản clone Times New Roman,
% đi kèm TeX Live, hỗ trợ đầy đủ dấu tiếng Việt.
\setmainfont{TeX Gyre Termes}
\setsansfont{TeX Gyre Heros}

% Chữ Hán: ưu tiên Source Han Serif SC (思源宋体) — chuẩn học thuật Trung Quốc.
% Nếu máy chưa cài, tự động dùng Songti SC (宋体, có sẵn trên macOS).
% Cả hai đều thuộc họ 宋体 nên diện mạo tài liệu không đổi.
\IfFontExistsTF{Source Han Serif SC}{
  \newcommand{\cjkfamilyname}{Source Han Serif SC}
}{
  \IfFontExistsTF{Songti SC}{
    \newcommand{\cjkfamilyname}{Songti SC}
  }{
    \newcommand{\cjkfamilyname}{Noto Serif CJK SC}
  }
}

% xeCJK tự chuyển font khi gặp ký tự Hán, không cần bọc \cjk{} thủ công.
% Đây là điều kiện bắt buộc: font Latin (TeX Gyre Termes) không có glyph chữ Hán,
% nếu không khai báo xeCJK thì mọi chữ Hán viết trần sẽ bị nuốt khỏi bản in.
\usepackage{xeCJK}
\setCJKfamilyfont{caehan}{\cjkfamilyname}[Scale=0.95]
\setCJKmainfont{\cjkfamilyname}[Scale=0.95]
\setCJKsansfont{\cjkfamilyname}[Scale=0.95]
\setCJKmonofont{\cjkfamilyname}[Scale=0.95]

% \cjk{} ép dùng font Hán trong ngữ cảnh cần thiết (ví dụ trong tiêu đề in đậm
% hoặc trong bảng có định dạng riêng). Dùng \CJKfamily chứ không phải tên font,
% vì tên font chỉ là văn bản — sẽ in ra chữ thay vì đổi font.
\newcommand{\cjk}[1]{{\CJKfamily{caehan}#1}}

% ---------- TOÁN HỌC ----------
\usepackage{amsmath, amssymb, amsthm, mathtools}
\usepackage{unicode-math}
\setmathfont{TeX Gyre Termes Math}

% ---------- BỐ CỤC TRANG ----------
% Lề phải chừa vùng an toàn của letterhead:
%   banner trên chiếm 0–19mm  → top tối thiểu 2,2cm
%   banner dưới chiếm 281–297mm → bottom tối thiểu 2,2cm
% Lề trên 2,6cm (26mm) để tiêu đề đầu trang không đè banner (banner hết ở 19mm).
\usepackage[a4paper, top=2.6cm, bottom=2.4cm, left=2.5cm, right=2.0cm]{geometry}
\usepackage{fancyhdr}
\usepackage{titlesec}
\usepackage{enumitem}
\usepackage{graphicx}
\usepackage{booktabs, array, longtable, tabularx}
\usepackage{xcolor}
\usepackage[most]{tcolorbox}
\usepackage[colorlinks=true, linkcolor=cae_blue, citecolor=cae_blue, urlcolor=cae_blue]{hyperref}

% ---------- BẢNG MÀU THƯƠNG HIỆU ----------
\definecolor{cae_blue}{RGB}{0, 51, 102}
\definecolor{cae_gold}{RGB}{196, 149, 60}
\definecolor{box_bg}{RGB}{245, 247, 250}
\definecolor{box_line}{RGB}{0, 51, 102}

% ---------- TIÊU ĐỀ MỤC ----------
% Đệm đầu trang: banner letterhead chiếm 0–19mm, mà lề trên là 22mm.
% \section* đầu trang (sau \clearpage) sát lề nên phải chừa thêm, nếu không
% sẽ đè lên banner.
\titleformat{\section}
  {\normalfont\large\bfseries\sffamily\color{cae_blue}}{\thesection.}{0.6em}{}
\titleformat{\subsection}
  {\normalfont\normalsize\bfseries\sffamily\color{cae_blue!85!black}}{\thesubsection}{0.6em}{}
\titlespacing*{\section}{0pt}{14pt}{6pt}
\titlespacing*{\subsection}{0pt}{10pt}{4pt}

% Tiêu đề không đánh số ở đầu trang: \section* sau \clearpage nằm sát lề trên
% nên đè lên banner letterhead. Dùng \vspace* ngay TRƯỚC \section* không đủ,
% vì titlesec chèn khoảng cách riêng. Cách chắc chắn là đệm bằng \null.
\newcommand{\sectiondautrang}[1]{\section*{#1}}

% Dạng bài dùng \subsection*{} nên cần định dạng riêng, đậm và có gạch chân màu
\titleformat{\subsubsection}
  {\normalfont\normalsize\bfseries\sffamily\color{cae_blue!85!black}}{\thesubsubsection}{0.6em}{}

% ---------- LETTERHEAD (ảnh nền toàn trang) ----------
% Ảnh letterhead của CAE Shanghai: banner xanh trên/dưới kèm logo chìm giữa.
% Tỉ lệ ảnh 1414x2000 = 0.7070, khớp A4 (0.7071) nên chèn thẳng không méo.
% Đặt ở lớp nền dưới cùng để chữ không bị ảnh che.
% Dùng \AddToShipoutPictureBG (không có dấu *) để áp dụng cho MỌI trang.
\usepackage{eso-pic}
\AddToShipoutPictureBG{%
  \AtPageLowerLeft{%
    \includegraphics[width=\paperwidth,height=\paperheight]{images/letterhead.png}%
  }%
}

% ---------- ĐẦU / CHÂN TRANG ----------
% Letterhead đã có banner trên và dưới, nên bỏ dòng chạy để tránh chồng chữ.
% Số trang đặt trong vùng nội dung, sát trên banner dưới, lề phải.
\pagestyle{fancy}
\setlength{\headheight}{0pt}
\fancyhf{}
\renewcommand{\headrulewidth}{0pt}
\renewcommand{\footrulewidth}{0pt}
\fancyfoot[R]{\raisebox{0.9cm}{\small\color{cae_blue}\thepage}}

% ---------- MÔI TRƯỜNG NỘI DUNG ----------
% Trong tcolorbox, LaTeX đặt parindent = parskip = 0, nên các đoạn văn trong
% khung dính liền nhau. Đặt lại khoảng cách đoạn cho mọi khung.
\newcommand{\khoangcachdoan}{\setlength{\parskip}{0.5em}\setlength{\parindent}{0pt}}
\tcbset{before upper={\khoangcachdoan}}

\newtcolorbox{dinhnghia}[1][]{
  colback=box_bg, colframe=box_line, boxrule=0.6pt, arc=1pt,
  fonttitle=\bfseries, title={Định nghĩa #1}, breakable}
\newtcolorbox{dinhly}[1][]{
  colback=box_bg, colframe=cae_gold, boxrule=0.6pt, arc=1pt,
  fonttitle=\bfseries, title={Tính chất #1}, breakable}
\newtcolorbox{vidu}[1][]{
  colback=white, colframe=black!40, boxrule=0.4pt, arc=1pt,
  fonttitle=\bfseries, title={Ví dụ #1}, breakable}
\newtcolorbox{baitap}[1][]{
  colback=white, colframe=cae_blue!50, boxrule=0.4pt, arc=1pt,
  fonttitle=\bfseries, title={Câu #1}, breakable}

% Ô đáp án: viền xanh đậm hơn, dùng trong file đáp án
\newtcolorbox{dapan}[1][]{
  colback=box_bg, colframe=cae_blue, boxrule=0.6pt, arc=1pt,
  fonttitle=\bfseries, title={Câu #1}, breakable}

% ---------- TÊN CHỨNG MINH ----------
\renewcommand{\proofname}{\normalfont\bfseries Chứng minh}

% ---------- MÔN HỌC ----------
% Đổi được ở file .tex trước khi % =====================================================================
%  CAE SHANGHAI — GÓI LỆNH DÙNG CHUNG CHO TÀI LIỆU LUYỆN THI CSCA TOÁN
%  Biên dịch bằng XeLaTeX. Mọi file .tex trong buổi học \input file này.
% =====================================================================

% ---------- NGÔN NGỮ & PHÔNG CHỮ ----------
\usepackage{fontspec}
\usepackage{polyglossia}
\setmainlanguage{vietnamese}
\setotherlanguage{english}

% Chữ Latin (Anh + Việt): TeX Gyre Termes — bản clone Times New Roman,
% đi kèm TeX Live, hỗ trợ đầy đủ dấu tiếng Việt.
\setmainfont{TeX Gyre Termes}
\setsansfont{TeX Gyre Heros}

% Chữ Hán: ưu tiên Source Han Serif SC (思源宋体) — chuẩn học thuật Trung Quốc.
% Nếu máy chưa cài, tự động dùng Songti SC (宋体, có sẵn trên macOS).
% Cả hai đều thuộc họ 宋体 nên diện mạo tài liệu không đổi.
\IfFontExistsTF{Source Han Serif SC}{
  \newcommand{\cjkfamilyname}{Source Han Serif SC}
}{
  \IfFontExistsTF{Songti SC}{
    \newcommand{\cjkfamilyname}{Songti SC}
  }{
    \newcommand{\cjkfamilyname}{Noto Serif CJK SC}
  }
}

% xeCJK tự chuyển font khi gặp ký tự Hán, không cần bọc \cjk{} thủ công.
% Đây là điều kiện bắt buộc: font Latin (TeX Gyre Termes) không có glyph chữ Hán,
% nếu không khai báo xeCJK thì mọi chữ Hán viết trần sẽ bị nuốt khỏi bản in.
\usepackage{xeCJK}
\setCJKfamilyfont{caehan}{\cjkfamilyname}[Scale=0.95]
\setCJKmainfont{\cjkfamilyname}[Scale=0.95]
\setCJKsansfont{\cjkfamilyname}[Scale=0.95]
\setCJKmonofont{\cjkfamilyname}[Scale=0.95]

% \cjk{} ép dùng font Hán trong ngữ cảnh cần thiết (ví dụ trong tiêu đề in đậm
% hoặc trong bảng có định dạng riêng). Dùng \CJKfamily chứ không phải tên font,
% vì tên font chỉ là văn bản — sẽ in ra chữ thay vì đổi font.
\newcommand{\cjk}[1]{{\CJKfamily{caehan}#1}}

% ---------- TOÁN HỌC ----------
\usepackage{amsmath, amssymb, amsthm, mathtools}
\usepackage{unicode-math}
\setmathfont{TeX Gyre Termes Math}

% ---------- BỐ CỤC TRANG ----------
% Lề phải chừa vùng an toàn của letterhead:
%   banner trên chiếm 0–19mm  → top tối thiểu 2,2cm
%   banner dưới chiếm 281–297mm → bottom tối thiểu 2,2cm
% Lề trên 2,6cm (26mm) để tiêu đề đầu trang không đè banner (banner hết ở 19mm).
\usepackage[a4paper, top=2.6cm, bottom=2.4cm, left=2.5cm, right=2.0cm]{geometry}
\usepackage{fancyhdr}
\usepackage{titlesec}
\usepackage{enumitem}
\usepackage{graphicx}
\usepackage{booktabs, array, longtable, tabularx}
\usepackage{xcolor}
\usepackage[most]{tcolorbox}
\usepackage[colorlinks=true, linkcolor=cae_blue, citecolor=cae_blue, urlcolor=cae_blue]{hyperref}

% ---------- BẢNG MÀU THƯƠNG HIỆU ----------
\definecolor{cae_blue}{RGB}{0, 51, 102}
\definecolor{cae_gold}{RGB}{196, 149, 60}
\definecolor{box_bg}{RGB}{245, 247, 250}
\definecolor{box_line}{RGB}{0, 51, 102}

% ---------- TIÊU ĐỀ MỤC ----------
% Đệm đầu trang: banner letterhead chiếm 0–19mm, mà lề trên là 22mm.
% \section* đầu trang (sau \clearpage) sát lề nên phải chừa thêm, nếu không
% sẽ đè lên banner.
\titleformat{\section}
  {\normalfont\large\bfseries\sffamily\color{cae_blue}}{\thesection.}{0.6em}{}
\titleformat{\subsection}
  {\normalfont\normalsize\bfseries\sffamily\color{cae_blue!85!black}}{\thesubsection}{0.6em}{}
\titlespacing*{\section}{0pt}{14pt}{6pt}
\titlespacing*{\subsection}{0pt}{10pt}{4pt}

% Tiêu đề không đánh số ở đầu trang: \section* sau \clearpage nằm sát lề trên
% nên đè lên banner letterhead. Dùng \vspace* ngay TRƯỚC \section* không đủ,
% vì titlesec chèn khoảng cách riêng. Cách chắc chắn là đệm bằng \null.
\newcommand{\sectiondautrang}[1]{\section*{#1}}

% Dạng bài dùng \subsection*{} nên cần định dạng riêng, đậm và có gạch chân màu
\titleformat{\subsubsection}
  {\normalfont\normalsize\bfseries\sffamily\color{cae_blue!85!black}}{\thesubsubsection}{0.6em}{}

% ---------- LETTERHEAD (ảnh nền toàn trang) ----------
% Ảnh letterhead của CAE Shanghai: banner xanh trên/dưới kèm logo chìm giữa.
% Tỉ lệ ảnh 1414x2000 = 0.7070, khớp A4 (0.7071) nên chèn thẳng không méo.
% Đặt ở lớp nền dưới cùng để chữ không bị ảnh che.
% Dùng \AddToShipoutPictureBG (không có dấu *) để áp dụng cho MỌI trang.
\usepackage{eso-pic}
\AddToShipoutPictureBG{%
  \AtPageLowerLeft{%
    \includegraphics[width=\paperwidth,height=\paperheight]{images/letterhead.png}%
  }%
}

% ---------- ĐẦU / CHÂN TRANG ----------
% Letterhead đã có banner trên và dưới, nên bỏ dòng chạy để tránh chồng chữ.
% Số trang đặt trong vùng nội dung, sát trên banner dưới, lề phải.
\pagestyle{fancy}
\setlength{\headheight}{0pt}
\fancyhf{}
\renewcommand{\headrulewidth}{0pt}
\renewcommand{\footrulewidth}{0pt}
\fancyfoot[R]{\raisebox{0.9cm}{\small\color{cae_blue}\thepage}}

% ---------- MÔI TRƯỜNG NỘI DUNG ----------
% Trong tcolorbox, LaTeX đặt parindent = parskip = 0, nên các đoạn văn trong
% khung dính liền nhau. Đặt lại khoảng cách đoạn cho mọi khung.
\newcommand{\khoangcachdoan}{\setlength{\parskip}{0.5em}\setlength{\parindent}{0pt}}
\tcbset{before upper={\khoangcachdoan}}

\newtcolorbox{dinhnghia}[1][]{
  colback=box_bg, colframe=box_line, boxrule=0.6pt, arc=1pt,
  fonttitle=\bfseries, title={Định nghĩa #1}, breakable}
\newtcolorbox{dinhly}[1][]{
  colback=box_bg, colframe=cae_gold, boxrule=0.6pt, arc=1pt,
  fonttitle=\bfseries, title={Tính chất #1}, breakable}
\newtcolorbox{vidu}[1][]{
  colback=white, colframe=black!40, boxrule=0.4pt, arc=1pt,
  fonttitle=\bfseries, title={Ví dụ #1}, breakable}
\newtcolorbox{baitap}[1][]{
  colback=white, colframe=cae_blue!50, boxrule=0.4pt, arc=1pt,
  fonttitle=\bfseries, title={Câu #1}, breakable}

% Ô đáp án: viền xanh đậm hơn, dùng trong file đáp án
\newtcolorbox{dapan}[1][]{
  colback=box_bg, colframe=cae_blue, boxrule=0.6pt, arc=1pt,
  fonttitle=\bfseries, title={Câu #1}, breakable}

% ---------- TÊN CHỨNG MINH ----------
\renewcommand{\proofname}{\normalfont\bfseries Chứng minh}

% ---------- MÔN HỌC ----------
% Đổi được ở file .tex trước khi \input{cae-style} bằng \renewcommand{\monhoc}{...}
\providecommand{\monhoc}{Môn Toán}

% ---------- TRANG BÌA GỌN ----------
% #1 = loại tài liệu (TÀI LIỆU CHUẨN BỊ, TÀI LIỆU BUỔI HỌC, ...)
% #2 = số buổi
% #3 = tên buổi tiếng Việt
% #4 = tên buổi tiếng Hán
% #5 = module và tỉ trọng
\newcommand{\trangbia}[5]{%
  \begin{center}
    \vspace*{0.6cm}
    {\large\bfseries\color{cae_blue} CAE SHANGHAI}\\[0.15cm]
    {\footnotesize Tài liệu luyện thi CSCA · \monhoc}\\[0.5cm]
    {\color{cae_gold}\rule{0.55\linewidth}{0.8pt}}\\[0.35cm]
    {\small\bfseries\color{cae_blue!80!black} #1}\\[0.35cm]
    {\Large\bfseries\color{cae_blue} BUỔI #2 — #3}\\[0.2cm]
    {\large\color{cae_blue!85!black} \cjk{#4}}\\[0.25cm]
    \ifx\relax#5\relax\else{\small #5}\\[0.35cm]\fi
    {\color{cae_gold}\rule{0.55\linewidth}{0.8pt}}\\[0.3cm]
    {\footnotesize Biên soạn: Hoàng Dương, CAE Shanghai \quad·\quad \today}
  \end{center}
  \vspace{0.3cm}
}

% ---------- TRANG BÌA CHO ĐỀ KIỂM TRA ----------
% Khác \trangbia ở chỗ không có chữ "BUỔI N" — dùng cho đề thi, không phải bài giảng.
% #1 = loại tài liệu, #2 = tên kỳ thi (Việt), #3 = tên kỳ thi (Hán), #4 = thông số
\newcommand{\biade}[4]{%
  \begin{center}
    \vspace*{0.6cm}
    {\large\bfseries\color{cae_blue} CAE SHANGHAI}\\[0.15cm]
    {\footnotesize Tài liệu luyện thi CSCA · \monhoc}\\[0.5cm]
    {\color{cae_gold}\rule{0.55\linewidth}{0.8pt}}\\[0.35cm]
    {\small\bfseries\color{cae_blue!80!black} #1}\\[0.35cm]
    {\Large\bfseries\color{cae_blue} #2}\\[0.2cm]
    {\large\color{cae_blue!85!black} \cjk{#3}}\\[0.25cm]
    {\small #4}\\[0.35cm]
    {\color{cae_gold}\rule{0.55\linewidth}{0.8pt}}\\[0.3cm]
    {\footnotesize Biên soạn: Hoàng Dương, CAE Shanghai \quad·\quad \today}
  \end{center}
  \vspace{0.3cm}
}

% ---------- TIỆN ÍCH ----------
% Nhãn từ khóa in đậm ở đầu đoạn, dùng cho Từ vựng / Từ khóa nhận dạng / Chú ý
% Khoảng cách đoạn do khung tự đặt (xem \khoangcachdoan ở trên),
% ở đây chỉ cần tách đoạn và bỏ thụt đầu dòng.
\newcommand{\nhan}[1]{\par\noindent\textbf{#1.}\ \ignorespaces}

% Bảng gọn, không đánh số
\newcommand{\bang}[1]{%
  \begin{center}#1\end{center}}

% Đánh dấu ô trống cho học sinh điền
\newcommand{\chotrong}{\makebox[3.2cm]{\dotfill}}
 bằng \renewcommand{\monhoc}{...}
\providecommand{\monhoc}{Môn Toán}

% ---------- TRANG BÌA GỌN ----------
% #1 = loại tài liệu (TÀI LIỆU CHUẨN BỊ, TÀI LIỆU BUỔI HỌC, ...)
% #2 = số buổi
% #3 = tên buổi tiếng Việt
% #4 = tên buổi tiếng Hán
% #5 = module và tỉ trọng
\newcommand{\trangbia}[5]{%
  \begin{center}
    \vspace*{0.6cm}
    {\large\bfseries\color{cae_blue} CAE SHANGHAI}\\[0.15cm]
    {\footnotesize Tài liệu luyện thi CSCA · \monhoc}\\[0.5cm]
    {\color{cae_gold}\rule{0.55\linewidth}{0.8pt}}\\[0.35cm]
    {\small\bfseries\color{cae_blue!80!black} #1}\\[0.35cm]
    {\Large\bfseries\color{cae_blue} BUỔI #2 — #3}\\[0.2cm]
    {\large\color{cae_blue!85!black} \cjk{#4}}\\[0.25cm]
    \ifx\relax#5\relax\else{\small #5}\\[0.35cm]\fi
    {\color{cae_gold}\rule{0.55\linewidth}{0.8pt}}\\[0.3cm]
    {\footnotesize Biên soạn: Hoàng Dương, CAE Shanghai \quad·\quad \today}
  \end{center}
  \vspace{0.3cm}
}

% ---------- TRANG BÌA CHO ĐỀ KIỂM TRA ----------
% Khác \trangbia ở chỗ không có chữ "BUỔI N" — dùng cho đề thi, không phải bài giảng.
% #1 = loại tài liệu, #2 = tên kỳ thi (Việt), #3 = tên kỳ thi (Hán), #4 = thông số
\newcommand{\biade}[4]{%
  \begin{center}
    \vspace*{0.6cm}
    {\large\bfseries\color{cae_blue} CAE SHANGHAI}\\[0.15cm]
    {\footnotesize Tài liệu luyện thi CSCA · \monhoc}\\[0.5cm]
    {\color{cae_gold}\rule{0.55\linewidth}{0.8pt}}\\[0.35cm]
    {\small\bfseries\color{cae_blue!80!black} #1}\\[0.35cm]
    {\Large\bfseries\color{cae_blue} #2}\\[0.2cm]
    {\large\color{cae_blue!85!black} \cjk{#3}}\\[0.25cm]
    {\small #4}\\[0.35cm]
    {\color{cae_gold}\rule{0.55\linewidth}{0.8pt}}\\[0.3cm]
    {\footnotesize Biên soạn: Hoàng Dương, CAE Shanghai \quad·\quad \today}
  \end{center}
  \vspace{0.3cm}
}

% ---------- TIỆN ÍCH ----------
% Nhãn từ khóa in đậm ở đầu đoạn, dùng cho Từ vựng / Từ khóa nhận dạng / Chú ý
% Khoảng cách đoạn do khung tự đặt (xem \khoangcachdoan ở trên),
% ở đây chỉ cần tách đoạn và bỏ thụt đầu dòng.
\newcommand{\nhan}[1]{\par\noindent\textbf{#1.}\ \ignorespaces}

% Bảng gọn, không đánh số
\newcommand{\bang}[1]{%
  \begin{center}#1\end{center}}

% Đánh dấu ô trống cho học sinh điền
\newcommand{\chotrong}{\makebox[3.2cm]{\dotfill}}
 bằng \renewcommand{\monhoc}{...}
\providecommand{\monhoc}{Môn Toán}

% ---------- TRANG BÌA GỌN ----------
% #1 = loại tài liệu (TÀI LIỆU CHUẨN BỊ, TÀI LIỆU BUỔI HỌC, ...)
% #2 = số buổi
% #3 = tên buổi tiếng Việt
% #4 = tên buổi tiếng Hán
% #5 = module và tỉ trọng
\newcommand{\trangbia}[5]{%
  \begin{center}
    \vspace*{0.6cm}
    {\large\bfseries\color{cae_blue} CAE SHANGHAI}\\[0.15cm]
    {\footnotesize Tài liệu luyện thi CSCA · \monhoc}\\[0.5cm]
    {\color{cae_gold}\rule{0.55\linewidth}{0.8pt}}\\[0.35cm]
    {\small\bfseries\color{cae_blue!80!black} #1}\\[0.35cm]
    {\Large\bfseries\color{cae_blue} BUỔI #2 — #3}\\[0.2cm]
    {\large\color{cae_blue!85!black} \cjk{#4}}\\[0.25cm]
    \ifx\relax#5\relax\else{\small #5}\\[0.35cm]\fi
    {\color{cae_gold}\rule{0.55\linewidth}{0.8pt}}\\[0.3cm]
    {\footnotesize Biên soạn: Hoàng Dương, CAE Shanghai \quad·\quad \today}
  \end{center}
  \vspace{0.3cm}
}

% ---------- TRANG BÌA CHO ĐỀ KIỂM TRA ----------
% Khác \trangbia ở chỗ không có chữ "BUỔI N" — dùng cho đề thi, không phải bài giảng.
% #1 = loại tài liệu, #2 = tên kỳ thi (Việt), #3 = tên kỳ thi (Hán), #4 = thông số
\newcommand{\biade}[4]{%
  \begin{center}
    \vspace*{0.6cm}
    {\large\bfseries\color{cae_blue} CAE SHANGHAI}\\[0.15cm]
    {\footnotesize Tài liệu luyện thi CSCA · \monhoc}\\[0.5cm]
    {\color{cae_gold}\rule{0.55\linewidth}{0.8pt}}\\[0.35cm]
    {\small\bfseries\color{cae_blue!80!black} #1}\\[0.35cm]
    {\Large\bfseries\color{cae_blue} #2}\\[0.2cm]
    {\large\color{cae_blue!85!black} \cjk{#3}}\\[0.25cm]
    {\small #4}\\[0.35cm]
    {\color{cae_gold}\rule{0.55\linewidth}{0.8pt}}\\[0.3cm]
    {\footnotesize Biên soạn: Hoàng Dương, CAE Shanghai \quad·\quad \today}
  \end{center}
  \vspace{0.3cm}
}

% ---------- TIỆN ÍCH ----------
% Nhãn từ khóa in đậm ở đầu đoạn, dùng cho Từ vựng / Từ khóa nhận dạng / Chú ý
% Khoảng cách đoạn do khung tự đặt (xem \khoangcachdoan ở trên),
% ở đây chỉ cần tách đoạn và bỏ thụt đầu dòng.
\newcommand{\nhan}[1]{\par\noindent\textbf{#1.}\ \ignorespaces}

% Bảng gọn, không đánh số
\newcommand{\bang}[1]{%
  \begin{center}#1\end{center}}

% Đánh dấu ô trống cho học sinh điền
\newcommand{\chotrong}{\makebox[3.2cm]{\dotfill}}


\begin{document}

\trangbia{KẾ HOẠCH GIẢNG DẠY}{13}{Vector, số phức và ôn tập hình học}{几何与代数综合}{}

\section{PHÂN BỔ THỜI GIAN}

\bang{
\begin{tabularx}{\linewidth}{@{}Xcc@{}}
\toprule
\textbf{Hoạt động} & \textbf{Tỉ trọng} & \textbf{Thời lượng} \\
\midrule
Tổng quan kết quả, phân loại lỗi & 10\% & 8 phút \\
Chữa nhanh câu cả lớp đúng: khái quát đáp án và cách nhận dạng & 25\% & 20 phút \\
Chữa kỹ câu nhiều học sinh sai & 45\% & 36 phút \\
Hệ thống nhóm Hình học và Đại số, kỹ thuật nhanh & 20\% & 16 phút \\
\bottomrule
\end{tabularx}}

Đề gồm 48 câu, học sinh làm trước buổi trong 60 phút. Khoảng 32 câu cả lớp làm đúng chỉ khái quát đáp án và từ khóa nhận dạng; 16 câu nhiều học sinh sai được chữa kỹ. Phần phân loại đúng/sai căn cứ trên bài nộp qua hệ thống, giáo viên chốt lại trước buổi.

\section{PHÂN BỐ CÂU THEO NHÓM NỘI DUNG}

\bang{
\begin{tabularx}{\linewidth}{@{}llcX@{}}
\toprule
\textbf{Nhóm} & \textbf{Câu} & \textbf{Số câu} & \textbf{Nguồn} \\
\midrule
Vector và số phức & 1–12 & 12 & kho đề CD7, trừ 3 câu đã dùng ở buổi 12 \\
Conic & 13–25 & 13 & kho đề CD6 và tự soạn \\
Đường thẳng và đường tròn & 26–36 & 11 & tự soạn \\
Hình học phẳng và không gian & 37–48 & 12 & kho đề CD5 và tự soạn \\
\textbf{Tổng} &  & \textbf{48} &  \\
\bottomrule
\end{tabularx}}

Vector và số phức chiếm 12 câu, cao hơn tỉ trọng thật (2–4 câu trong đề CSCA). Phần dôi ra là chủ ý: đây là chuyên đề vừa học ở buổi 12 và là nội dung học sinh yếu nhất, cần luyện dày trước Kiểm tra giữa kỳ 2. Conic 13 câu phản ánh đúng trọng số lớn nhất của bài kiểm tra đó.

\subsection{Phân bố đáp án}

\bang{
\begin{tabularx}{\linewidth}{@{}Xcccc@{}}
\toprule
\textbf{Nhóm} & \textbf{A} & \textbf{B} & \textbf{C} & \textbf{D} \\
\midrule
Câu 1–12 & 5 & 1 & 3 & 3 \\
Câu 13–25 & 2 & 3 & 3 & 5 \\
Câu 26–36 & 3 & 5 & 2 & 1 \\
Câu 37–48 & 2 & 3 & 4 & 3 \\
\textbf{Tổng} & \textbf{12} & \textbf{12} & \textbf{12} & \textbf{12} \\
\bottomrule
\end{tabularx}}

Bảng đáp án cân đều bốn phương án trên toàn đề, nhưng từng phần có độ lệch cục bộ (phần I chỉ có 1 câu B, phần II có 5 câu D). Đây là điều bình thường ở đề 48 câu; nhắc học sinh rằng phân bố đều là đặc điểm của đề chuẩn, không phải gợi ý để đoán.

\section{LỖI HAY MẮC THEO NHÓM}

\subsection{Vector và số phức}

\bang{
\begin{tabularx}{\linewidth}{@{}llX@{}}
\toprule
\textbf{Câu} & \textbf{Lỗi} & \textbf{Cách cảnh báo} \\
\midrule
3 & Quên đổi dấu thành phần giữa khi tính tích có hướng & Viết công thức khai triển định thức lên bảng, khoanh tròn dấu trừ ở số hạng giữa \\
4 & Nhầm $\overrightarrow{DB}$ với $\overrightarrow{BD}$ & Vẽ hình, chỉ rõ gốc và ngọn của vectơ trước khi biến đổi \\
8 & Chỉ lấy $k=2$, bỏ nghiệm $k=-2$ & Nhắc: phương trình $k^2=4$ luôn có hai nghiệm; đề không ràng buộc dấu thì giữ cả hai \\
10 & Chọn $z_1$ vì quên bước cộng $z_2$ & Yêu cầu gạch chân cụm từ chỉ phép toán cuối cùng trong đề \\
12 & Không biết dùng trung điểm để đưa về $\dfrac{1}{2}(\overrightarrow{AB}+\overrightarrow{AC})$ & Đây là công thức gốc của mọi bài trung tuyến, phải thuộc \\
\bottomrule
\end{tabularx}}

\subsection{Conic}

\bang{
\begin{tabularx}{\linewidth}{@{}llX@{}}
\toprule
\textbf{Câu} & \textbf{Lỗi} & \textbf{Cách cảnh báo} \\
\midrule
14 & Chỉ xét một hướng trục, bỏ trường hợp còn lại & Gặp tỉ số $\dfrac{b}{a}$ mà đề không nói trục nào thì luôn xét hai trường hợp \\
16 & Không nhận ra dấu bằng khi ba điểm thẳng hàng & Vẽ hình, chốt: bất đẳng thức tam giác chỉ chặt khi ba điểm không thẳng hàng \\
18 & Đọc $\dfrac{b}{a}$ thành $\dfrac{a}{b}$ & Viết to hai chữ cái lên bảng cạnh nhau, đọc thành tiếng từng lần \\
20 & Nhầm vị trí trục thực khi số hạng $y^2$ đứng trước & Quy tắc: nhìn mẫu của số hạng dương, đó là trục thực \\
22 & Mặc định tiêu điểm luôn trên $Ox$ & So sánh hai mẫu trước khi kết luận \\
\bottomrule
\end{tabularx}}

\subsection{Đường thẳng và đường tròn}

\bang{
\begin{tabularx}{\linewidth}{@{}llX@{}}
\toprule
\textbf{Câu} & \textbf{Lỗi} & \textbf{Cách cảnh báo} \\
\midrule
26 & Nhầm $\tan 120^\circ$ với $\tan 60^\circ$ & Góc tù cho hệ số góc âm, kiểm tra dấu trước khi chọn \\
31 & Đọc tọa độ tâm sai dấu & Viết mẫu: $D=-6$ thì $-\dfrac{D}{2}=3$, không phải $-3$ \\
32 & Lấy khoảng cách từ tâm tới đường thẳng, quên trừ bán kính & Vẽ hình: khoảng cách ngắn nhất là từ điểm trên đường tròn, luôn là $d-r$ \\
33 & So sánh sai $d$ với $r$ & Nhắc lại ba trường hợp, yêu cầu viết ra giấy trước khi chọn \\
35 & Tính $r^2$ rồi dừng, không lấy căn & Đề hỏi bán kính, không hỏi bình phương bán kính \\
\bottomrule
\end{tabularx}}

\subsection{Hình học phẳng và không gian}

\bang{
\begin{tabularx}{\linewidth}{@{}llX@{}}
\toprule
\textbf{Câu} & \textbf{Lỗi} & \textbf{Cách cảnh báo} \\
\midrule
38 & Lấy $\lvert MN\rvert$ làm bán kính, quên chia đôi & Đường kính gấp đôi bán kính; đọc kỹ đề hỏi đường kính hay bán kính \\
39 & Đổi dấu sai khi lấy đối xứng qua trục & Trục nào thì tọa độ ứng với trục đó giữ nguyên, hai tọa độ còn lại đổi dấu \\
41 & Chọn phương án giữ nguyên cả ba tọa độ & Phương án đó chính là điểm đã cho, không phải điểm đối xứng \\
44 & Nhầm bán kính với bình phương bán kính trong phương trình mặt cầu & Vế phải luôn là $R^2$, không phải $R$ \\
46 & Lẫn công thức diện tích mặt cầu với thể tích & Viết hai công thức cạnh nhau: $S=4\pi R^2$ và $V=\dfrac{4}{3}\pi R^3$ \\
\bottomrule
\end{tabularx}}

\section{DANH SÁCH CÂU CẦN CHỮA KỸ}

Mười sáu câu dưới đây chiếm phần lớn thời gian của khối chữa kỹ. Chọn theo mức độ học sinh hay sai và mức độ bao phủ kỹ thuật, không chữa hết mọi câu sai.

\subsection{Ưu tiên 1: bắt buộc chữa (8 câu, khoảng 20 phút)}

\bang{
\begin{tabularx}{\linewidth}{@{}llXc@{}}
\toprule
\textbf{Câu} & \textbf{Mã câu} & \textbf{Lý do} & \textbf{Thời lượng} \\
\midrule
8 & \texttt{CAE-M-CD7-08.4-H-016} & Bẫy nghiệm kép của phương trình song song, học sinh hay bỏ nghiệm âm & 2,5 phút \\
12 & \texttt{CAE-M-CD7-LX2-H-039} & Kỹ thuật trung điểm cộng tích vô hướng, dạng câu khó nhất phần vector & 3 phút \\
14 & \texttt{CAE-M-CD6-LX2-E-029} & Bẫy hai hướng trục của hypebol & 2,5 phút \\
16 & \texttt{CAE-M-CD6-LX3-H-033} & Bất đẳng thức tam giác cộng định nghĩa hypebol, câu khó nhất đề & 3,5 phút \\
18 & tự soạn & Nhầm $\dfrac{b}{a}$ với $\dfrac{a}{b}$, lỗi lặp lại ở mọi đề & 2 phút \\
32 & tự soạn & Quên trừ bán kính khi tìm khoảng cách nhỏ nhất & 2 phút \\
38 & \texttt{CAE-M-CD5-08.8-H-044} & Mặt cầu đường kính $MN$: quên chia đôi & 2,5 phút \\
39 & \texttt{CAE-M-CD5-LX1-E-047} & Quy tắc dấu đối xứng qua trục trong không gian & 2 phút \\
\bottomrule
\end{tabularx}}

\subsection{Ưu tiên 2: chữa nếu còn thời gian (8 câu, khoảng 16 phút)}

\bang{
\begin{tabularx}{\linewidth}{@{}llXc@{}}
\toprule
\textbf{Câu} & \textbf{Mã câu} & \textbf{Lý do} & \textbf{Thời lượng} \\
\midrule
3 & \texttt{CAE-M-CD7-08.3-M-007} & Tích có hướng, kỹ thuật cơ bản nhưng dễ sai dấu & 2 phút \\
10 & \texttt{CAE-M-CD7-08.7-M-024} & Đọc đề thiếu bước tính cuối & 1,5 phút \\
20 & tự soạn & Xác định trục thực của hypebol theo số hạng dương & 2 phút \\
22 & tự soạn & Tiêu điểm elip khi trục lớn trên $Oy$ & 2 phút \\
26 & tự soạn & Hệ số góc của góc nghiêng tù & 1,5 phút \\
31 & tự soạn & Tọa độ tâm đường tròn, lỗi dấu & 2 phút \\
35 & tự soạn & Phân biệt bán kính và bình phương bán kính & 1,5 phút \\
44 & tự soạn & Phương trình mặt cầu, vế phải là $R^2$ & 2 phút \\
\bottomrule
\end{tabularx}}

\subsection{Thứ tự ưu tiên khi hết giờ}

Nếu chỉ còn 10 phút cuối, chữa bốn câu theo thứ tự: \textbf{Câu 16, Câu 12, Câu 8, Câu 38}. Bốn câu này đại diện cho bốn kỹ thuật khác nhau và là những lỗi học sinh sẽ gặp lại trong Kiểm tra giữa kỳ 2.

\section{KHỐI 1 — TỔNG QUAN VÀ PHÂN LOẠI LỖI (8 phút)}

\subsection{Mục tiêu}

Học sinh biết mình sai ở nhóm nào và vì sao, trước khi nghe chữa từng câu.

\subsection{Kịch bản}

\textbf{Phút 1–3 · Chiếu phổ điểm.} Nêu điểm trung bình, điểm cao nhất, điểm thấp nhất. Không đọc tên. Nhấn mạnh khoảng cách giữa nhóm làm tốt và nhóm làm chậm nằm ở tốc độ đọc đề, không phải ở kiến thức.

\textbf{Phút 4–6 · Ba nhóm lỗi chính.} Chiếu bảng phân loại lỗi, chốt ba nguyên nhân chiếm phần lớn số câu sai:

1. \textbf{Bỏ sót trường hợp.} Câu 8 và Câu 14 là ví dụ: đề không ràng buộc dấu hoặc không nói trục, học sinh chỉ lấy một nghiệm.
2. \textbf{Nhầm công thức cùng họ.} Tiệm cận hypebol $\dfrac{b}{a}$ với $\dfrac{a}{b}$, diện tích mặt cầu với thể tích, bán kính với bình phương bán kính.
3. \textbf{Quên bước cuối.} Câu 10 quên cộng $z_2$, Câu 32 quên trừ bán kính, Câu 38 quên chia đôi.

\textbf{Phút 7–8 · Đối chiếu tỉ lệ.} Chiếu bảng phân bố câu theo nhóm và nói rõ: phần Conic và Hình học giải tích chiếm 24 trên 48 câu. Đây là lý do hai nhóm này được chữa kỹ nhất, và cũng là trọng tâm của bài kiểm tra sắp tới.

\section{KHỐI 2 — CHỮA NHANH CÂU CẢ LỚP ĐÚNG (20 phút)}

\subsection{Mục tiêu}

Học sinh hệ thống lại cách nhận dạng, không phải nghe lại lời giải. Với mỗi câu, chỉ nêu \textbf{từ khóa nhận dạng} và \textbf{công thức áp dụng}, không biến đổi chi tiết.

\subsection{Kịch bản}

Ba mươi hai câu còn lại chữa theo lô, mỗi lô 4–6 câu, khoảng 3 phút mỗi lô. Cách làm: đọc từ khóa tiếng Trung, hỏi cả lớp đây là dạng gì, chốt công thức, đọc đáp án.

\bang{
\begin{tabularx}{\linewidth}{@{}lllX@{}}
\toprule
\textbf{Lô} & \textbf{Câu} & \textbf{Từ khóa nhận dạng} & \textbf{Công thức chốt} \\
\midrule
1 & 1, 2, 4, 5, 7 & 垂直, 纯虚数, 向量减法, 方向向量, 平行 & $a\cdot b=0$; phần thực bằng 0; $\overrightarrow{DB}=\overrightarrow{AB}-\overrightarrow{AD}$; $\cos\theta=\dfrac{\lvert d_1\cdot d_2\rvert}{\lvert d_1\rvert\lvert d_2\rvert}$; tỉ số thành phần không đổi \\
2 & 6, 9, 11 & 复数除法, 夹角相等, 投影 & Nhân liên hợp mẫu; hai mẫu bằng nhau thì đồng nhất tử \\
3 & 13, 15, 17, 19, 21 & 抛物线焦点, 准线, 离心率, 长轴长 & Bán kính qua tiêu $x+\dfrac{p}{2}$; $e=\dfrac{c}{a}$; trục lớn $2a$ \\
4 & 23, 24, 25 & 准线方程, 定义, 焦半径 & Đường chuẩn khác phía hướng mở; $\lvert PF_1\rvert+\lvert PF_2\rvert=2a$ \\
5 & 27, 28, 29, 30 & 平行, 交点, 倾斜角 & Song song giữ hệ số góc; giải hệ hai phương trình; $k=\tan\alpha$ \\
6 & 33, 34, 36, 37 & 位置关系, 圆心, 点到直线距离 & So $d$ với $r$; tâm $\left(-\dfrac{D}{2},-\dfrac{E}{2}\right)$; $d=\dfrac{\lvert Ax_0+By_0+C\rvert}{\sqrt{A^2+B^2}}$ \\
7 & 40, 41, 42, 43, 45, 46, 47, 48 & 坐标平面, 对称点, 中点, 侧面积, 表面积, 体对角线 & Cao độ bằng 0 thuộc $xOy$; $S_{xq}=2\pi rl$; $S=4\pi R^2$; $d=\sqrt{a^2+b^2+c^2}$ \\
\bottomrule
\end{tabularx}}

\subsection{Điểm cần chốt}

Sau khối này, học sinh phải nhận ra ngay ba cặp từ khóa dễ lẫn:

\begin{itemize}[leftmargin=1.4em, itemsep=2pt]
\item 渐近线 (tiệm cận) so với 准线 (đường chuẩn)
\item 长轴 (trục lớn) so với 短轴 (trục nhỏ)
\item 侧面积 (diện tích xung quanh) so với 表面积 (diện tích toàn phần)
\end{itemize}

\section{KHỐI 3 — CHỮA KỸ CÂU NHIỀU HỌC SINH SAI (36 phút)}

\subsection{Mục tiêu}

Học sinh hiểu \textbf{nguyên nhân sai}, không chỉ biết đáp án đúng.

\subsection{Cấu trúc mỗi câu (khoảng 2–3,5 phút)}

1. \textbf{Đọc đề tiếng Trung} (20 giây): giáo viên đọc, học sinh chỉ ra từ khóa.
2. \textbf{Hỏi lớp} (20 giây): "Ai chọn phương án nào, vì sao?"
3. \textbf{Chỉ ra chỗ sai} (40 giây): gọi học sinh chọn sai giải thích, giáo viên khoanh vùng lỗi.
4. \textbf{Giảng cách đúng} (60–90 giây): trình bày theo thứ tự nhận dạng, công thức, biến đổi ngắn.
5. \textbf{Chốt quy tắc} (20 giây): một câu để học sinh ghi vào vở.

\subsection{Kịch bản theo nhóm}

\textbf{Phút 1–10 · Bốn câu vector và số phức (Câu 3, 8, 10, 12).} Đây là nhóm học sinh vừa học ở buổi 12 nên còn nhớ lý thuyết, sai chủ yếu do bất cẩn. Dành nhiều thời gian nhất cho Câu 12: vẽ tam giác, chỉ rõ $D$ là trung điểm nên $\overrightarrow{AD}=\dfrac{1}{2}(\overrightarrow{AB}+\overrightarrow{AC})$, sau đó mới thay biểu thức theo $a$ và $b$. Nhấn mạnh bước bình phương để đưa về tích vô hướng.

\textbf{Phút 11–24 · Bốn câu conic (Câu 14, 16, 18, 20).} Với Câu 14, yêu cầu học sinh tự viết cả hai phương trình ứng với hai hướng trục rồi thay điểm, để thấy cả hai đều thỏa mãn. Với Câu 16, vẽ hypebol và đường thẳng $PF_2$, chỉ rõ điểm $Q$ cần tìm là giao điểm gần $F_2$ nhất; đây là câu khó nhất đề, chấp nhận nếu lớp không theo kịp, nhưng phải chốt được ý tưởng dùng bất đẳng thức tam giác.

\textbf{Phút 25–31 · Ba câu đường thẳng và đường tròn (Câu 26, 31, 32).} Ba lỗi dấu và bước cuối, chữa nhanh. Vẽ hình minh họa cho Câu 32: khoảng cách nhỏ nhất từ điểm trên đường tròn tới đường thẳng nằm ngoài là $d-r$.

\textbf{Phút 32–36 · Năm câu hình học phẳng và không gian (Câu 35, 38, 39, 41, 44).} Chốt hai quy tắc dấu: đối xứng qua trục nào thì giữ tọa độ trục đó, và vế phải phương trình mặt cầu luôn là $R^2$.

\section{KHỐI 4 — HỆ THỐNG NHÓM HÌNH HỌC VÀ ĐẠI SỐ, KỸ THUẬT NHANH (16 phút)}

\subsection{Mục tiêu}

Học sinh có một bảng tra công thức dùng chung cho cả ba nhóm, và biết ba kỹ thuật cắt thời gian.

\subsection{Bảng hệ thống công thức}

\bang{
\begin{tabularx}{\linewidth}{@{}llX@{}}
\toprule
\textbf{Nhóm} & \textbf{Đại lượng} & \textbf{Công thức} \\
\midrule
Elip & Tiêu cự, tâm sai & $c^2=a^2-b^2$, $e=\dfrac{c}{a}$ với $a>b>0$ \\
Hypebol & Tiêu cự, tiệm cận & $c^2=a^2+b^2$, $y=\pm\dfrac{b}{a}x$ (trục thực trên $Ox$) \\
Parabol & Tiêu điểm, đường chuẩn, bán kính qua tiêu & $F\left(\dfrac{p}{2},0\right)$, $x=-\dfrac{p}{2}$, $\lvert PF\rvert=x_P+\dfrac{p}{2}$ \\
Đường tròn & Tâm, bán kính & $I\left(-\dfrac{D}{2},-\dfrac{E}{2}\right)$, $r=\sqrt{\dfrac{D^2+E^2}{4}-F}$ \\
Đường thẳng & Khoảng cách điểm tới đường thẳng & $d=\dfrac{\lvert Ax_0+By_0+C\rvert}{\sqrt{A^2+B^2}}$ \\
Mặt cầu & Phương trình & $(x-a)^2+(y-b)^2+(z-c)^2=R^2$ \\
Vector & Góc, tích có hướng & $\cos\theta=\dfrac{a\cdot b}{\lvert a\rvert\lvert b\rvert}$; $a\times b$ đổi dấu thành phần giữa \\
Số phức & Môđun, liên hợp & $\lvert z\rvert^2=z\overline{z}$ \\
\bottomrule
\end{tabularx}}

\subsection{Ba kỹ thuật cắt thời gian}

\textbf{Kỹ thuật 1: thay đáp án vào đề.} Dùng cho câu hỏi tìm tham số (Câu 8, Câu 29) và câu hỏi tìm phương trình đường (Câu 27, Câu 34). Thay một giá trị, kiểm tra một đẳng thức, loại được hai phương án trong khoảng 15 giây.

\textbf{Kỹ thuật 2: đọc từ khóa để chọn công thức trước khi đọc số.} Câu hỏi có 渐近线 thì nghĩ ngay $\dfrac{b}{a}$, có 准线 thì nghĩ ngay $\dfrac{p}{2}$, có 离心率 thì nghĩ ngay $\dfrac{c}{a}$. Nhận ra từ khóa trong 3 giây, không cần dịch cả đề.

\textbf{Kỹ thuật 3: ước lượng và loại trừ.} Với câu tính khoảng cách (Câu 32, Câu 36, Câu 37), ước lượng bậc độ lớn trước: khoảng cách từ điểm tới đường thẳng trong mặt phẳng thường nhỏ hơn $5$, trong không gian thường nhỏ hơn $10$. Phương án lệch hẳn bậc độ lớn bị loại ngay.

\subsection{Điểm cần chốt}

Ba câu viết lên góc bảng và để nguyên khi học sinh ra về:

1. \textbf{Đề không ràng buộc dấu thì giữ cả hai nghiệm.}
2. \textbf{Vế phải phương trình mặt cầu luôn là $R^2$.}
3. \textbf{Khoảng cách nhỏ nhất từ đường tròn tới đường thẳng là $d-r$, không phải $d$.}

\section{CHUẨN BỊ CHO KIỂM TRA GIỮA KỲ 2 (5 phút cuối, nằm trong khối 4)}

\subsection{Thông tin bài kiểm tra}

\bang{
\begin{tabularx}{\linewidth}{@{}lX@{}}
\toprule
\textbf{Nội dung} & \textbf{Thông tin} \\
\midrule
Thời điểm & Sau buổi 13, làm trực tuyến trên hệ thống có tính giờ \\
Số câu & 36 câu trắc nghiệm \\
Thời gian & 50 phút \\
Tỉ trọng điểm & 15\% điểm khóa học \\
Điều kiện & Không dùng máy tính \\
\bottomrule
\end{tabularx}}

\subsection{Ma trận phân bổ}

\bang{
\begin{tabularx}{\linewidth}{@{}Xcl@{}}
\toprule
\textbf{Nội dung} & \textbf{Số câu} & \textbf{Buổi tương ứng} \\
\midrule
Conic & 13 & Buổi 10 \\
Đường thẳng và đường tròn & 9 & Buổi 9 \\
Vector và số phức & 7 & Buổi 12 \\
Hình học phẳng và không gian & 7 & Buổi 8 \\
\bottomrule
\end{tabularx}}

Conic chiếm 13 trên 36 câu, tức hơn một phần ba bài kiểm tra. Đây là nhóm nội dung trọng tâm nhất và cũng là nhóm học sinh hay mất điểm nhất do nhầm công thức tiệm cận và vị trí trục.

\subsection{Ba việc học sinh phải làm trước khi kiểm tra}

1. \textbf{Học thuộc bảng công thức ba đường conic.} Elip dùng $c^2=a^2-b^2$, hypebol dùng $c^2=a^2+b^2$. Nhầm dấu ở đây làm sai cả câu.
2. \textbf{Luyện lại phần vector và số phức.} Đây là chuyên đề vừa học, dễ lấy điểm nhất trong bốn nhóm, nhưng chỉ có 7 câu nên mỗi câu sai mất khoảng 2 điểm phần trăm.
3. \textbf{Bấm giờ khi làm đề.} Bài kiểm tra có 50 phút cho 36 câu, tức khoảng 83 giây mỗi câu. Nhanh hơn đề thi thật một chút, nên không có chỗ cho biến đổi dài dòng.

\subsection{Dặn dò}

Không giao bài tập về nhà theo khối lượng thông thường. Nhiệm vụ của học sinh là \textbf{ôn lại đề này và làm lại những câu đã sai}, đặc biệt là 16 câu trong danh sách chữa kỹ. Ngoài ra ôn lại đề buổi 11 để nhớ phần hình học phẳng và không gian, vì nhóm này chiếm 7 câu trong bài kiểm tra.

Buổi 14 học về xác suất, nội dung độc lập với bài kiểm tra.

\section{GHI CHÚ SAU BUỔI HỌC}

\emph{(Điền sau khi dạy xong, dùng để điều chỉnh buổi 16 và các buổi lặp lại.)}

\bang{
\begin{tabularx}{\linewidth}{@{}Xl@{}}
\toprule
\textbf{Nội dung} & \textbf{Ghi nhận} \\
\midrule
Thời gian thực tế từng khối &  \\
Câu học sinh hỏi nhiều nhất &  \\
Nhóm nội dung yếu nhất &  \\
Điều chỉnh cho lần dạy sau &  \\
\bottomrule
\end{tabularx}}

\emph{Tài liệu nội bộ, CAE SHANGHAI.}
\vspace{0.8cm}
\begin{center}
{\footnotesize\color{cae_blue} Tài liệu nội bộ, CAE SHANGHAI}
\end{center}

\end{document}
